\documentclass[a4paper, 12pt]{article}
\usepackage{amsmath}
\usepackage{amssymb}
\usepackage{amsthm}
\usepackage{amsfonts}
\usepackage{mathtools}
\usepackage{hyperref}
\usepackage[margin=2.5cm]{geometry}

\renewcommand{\arraystretch}{1.5}
\renewcommand{\qedsymbol}{$\blacksquare$}

\newcommand{\dd}{\mathrm{d}}
\newcommand{\Hom}{\operatorname{Hom}}
\newcommand{\R}{\mathbb{R}}
\newcommand{\C}{\mathbb{C}}
\newcommand{\diag}{\operatorname{diag}}
\newcommand{\SE}{\operatorname{SE}(3)}
\newcommand{\SO}{\operatorname{SO}(3)}
\newcommand{\lmax}{\ell_{\text{max}}}

\newtheorem{theorem}{Theorem}[section]
\newtheorem{corollary}[theorem]{Corollary}
\newtheorem{lemma}[theorem]{Lemma}
\newtheorem{prop}[theorem]{Proposition}
\theoremstyle{definition}
\newtheorem{definition}[theorem]{Definition}
\newtheorem{example}[theorem]{Example}
\newtheorem{remark}[theorem]{Remark}

\title{Fast and Memory-Efficient Low-Rank $\operatorname{SE}(3)$-Equivariant Networks}
\author{Hamish M. Blair}

\hypersetup{
    colorlinks,
    citecolor=black,
    filecolor=black,
    linkcolor=black,
    urlcolor=black
}

\begin{document}

\maketitle

\section{Introduction}

Modern equivariant networks are predicated on the computation of equivariant weight matrices $W : \mathbb{R}^3 \to \Hom(\mathbb{R}^{2 \ell_1 + 1}, \mathbb{R}^{2 \ell_2 + 1})$, satisfying $W(g \cdot x) = D^{\ell_2}(g)^{\top} W(x) D^{\ell_1}(g)$ for all $x \in \mathbb{R}^3$ and $g \in \SO$. These are typically constructed as linear combinations of $2 \min(\ell_1, \ell_2) + 1$ basis elements
\begin{align*}
	W_{\ell_1 \ell_2}^{\ell}(x) = \sum_{m = - \ell}^{\ell} Q_{\ell_1 \ell_2}^{\ell m} Y_{\ell m}(x),
\end{align*}
where $Q_{\ell_1 \ell_2}^{\ell m}$ are Clebsch-Gordan coefficients. If we consider all basis elements up to a maximum degree $\lmax$, there are $\mathcal{O}(\lmax^3)$ distinct basis matrices, leading to algorithms of order $\mathcal{O}(\lmax^9)$ in time and $\mathcal{O}(\lmax^7)$ in space.

The central observation of this paper is that the space of equivariant matrices is \emph{diagonalized} by the Wigner-D matrices---in the sense of the real Schur's lemma, where invariant subspaces may be one- or two-dimensional. This yields a low-rank representation requiring only $\mathcal{O}(\lmax^2)$ basis vectors rather than $\mathcal{O}(\lmax^3)$ basis matrices, reducing complexity to $\mathcal{O}(\lmax^6)$ in time and $\mathcal{O}(\lmax^4)$ in space.


\section{Theoretical Results}

\subsection{Branching Rules}

\begin{lemma}\label{lem:branching}
Let $G = \SO$ and $H = \operatorname{SO}(2)$. The branching rule for the restriction of real irreducible representations is
\begin{align*}
	D^{\ell}|_{H} \cong \rho^0 \oplus \bigoplus_{m = 1}^{\ell} \rho^m,
\end{align*}
where $\rho^0$ is the trivial one-dimensional representation and $\rho^m$ for $m \geq 1$ is the two-dimensional representation given by rotation by angle $m\theta$.
\end{lemma}

\begin{proof}
The character of $D^{\ell}$ restricted to rotations by angle $\theta$ about a fixed axis is the Dirichlet kernel $\sum_{m=-\ell}^{\ell} e^{im\theta}$. In real form, this decomposes as $1 + \sum_{m=1}^{\ell} 2\cos(m\theta)$, which is precisely the character of $\rho^0 \oplus \bigoplus_{m=1}^{\ell} \rho^m$.
\end{proof}


\subsection{Diagonalization by Wigner-D Matrices}

The following theorem is the key structural result.

\begin{theorem}\label{thm:main}
Let $W : \mathbb{R}^3 \to \Hom(\mathbb{R}^{2\ell + 1}, \mathbb{R}^{2\ell + 1})$ be $\SO$-equivariant. For any nonzero $x \in \mathbb{R}^3$, let $g_x \in \SO$ be a rotation taking $e_z$ to $\hat{x} = x/\|x\|$. Then $W(x)$ is block-diagonalized by the Wigner-D matrix $D^{\ell}(g_x)$:
\begin{align}
	D^{\ell}(g_x)^{\top} W(x) D^{\ell}(g_x) = \begin{pmatrix} \lambda_0 & & & \\ & A_1 & & \\ & & \ddots & \\ & & & A_{\ell} \end{pmatrix},
\end{align}
where $\lambda_0 \in \R$ and each $A_m$ for $m \geq 1$ is a $2 \times 2$ block of the form
\begin{align}
	A_m = \begin{pmatrix} a_m & b_m \\ -b_m & a_m \end{pmatrix}
\end{align}
for some $a_m, b_m \in \R$ depending on $x$ and $W$.
\end{theorem}

\begin{proof}
Let $H = \operatorname{stab}_{\SO}(x) \cong \operatorname{SO}(2)$. For any $h \in H$, the equivariance condition gives
\begin{align*}
	W(x) = W(h \cdot x) = D^{\ell}(h)^{\top} W(x) D^{\ell}(h),
\end{align*}
so $W(x)$ commutes with $D^{\ell}(h)$ for all $h \in H$. Thus $W(x)$ is an $H$-equivariant endomorphism.

By the branching rule (Lemma \ref{lem:branching}), the space $\R^{2\ell+1}$ decomposes under $H$ as
\begin{align*}
	\R^{2\ell+1} = V_0 \oplus \bigoplus_{m=1}^{\ell} V_m,
\end{align*}
where $\dim V_0 = 1$ and $\dim V_m = 2$ for $m \geq 1$. The Wigner-D matrix $D^{\ell}(g_x)$ effects this decomposition: its $m = 0$ column spans $V_0$, while its $m$ and $-m$ columns (for $m > 0$) span $V_m$.

By the real Schur's lemma, $W(x)$ acts as a scalar on $V_0$ and as a rotation-scaling (complex-type map) on each $V_m$. Conjugating by $D^{\ell}(g_x)$ transforms $W(x)$ to this block-diagonal form.
\end{proof}

\begin{remark}
The block structure reflects the type of each irreducible $H$-representation:
\begin{itemize}
	\item The $m = 0$ subspace is of \emph{real type}: the intertwiner space is $\R$.
	\item Each $m > 0$ subspace is of \emph{complex type}: the intertwiner space is $\C \cong \R^2$, realized as rotation-scalings.
\end{itemize}
\end{remark}


\subsection{Consequences}

\begin{corollary}
If $W_1, W_2 : \mathbb{R}^3 \to \Hom(\mathbb{R}^{2\ell + 1}, \mathbb{R}^{2\ell + 1})$ are both $\SO$-equivariant, then $W_1(x)$ and $W_2(x)$ commute for any $x \in \mathbb{R}^3$.
\end{corollary}

\begin{proof}
Both are simultaneously block-diagonalized by $D^{\ell}(g_x)$, and maps of complex type on the same subspace commute.
\end{proof}

\begin{corollary}
For the case $\ell_1 \neq \ell_2$, the decomposition becomes
\begin{align*}
	D^{\ell_2}(g_x)^{\top} W(x) D^{\ell_1}(g_x) = \begin{pmatrix} \lambda_0 & & \\ & A_1 & \\ & & \ddots \end{pmatrix},
\end{align*}
where the block structure has $\min(\ell_1, \ell_2) + 1$ blocks (one $1 \times 1$ and $\min(\ell_1, \ell_2)$ blocks of size $2 \times 2$). This gives
\begin{align*}
	\dim \Hom_{\SO}\bigl(\R^3, \Hom(\R^{2\ell_1+1}, \R^{2\ell_2+1})\bigr) = 2\min(\ell_1, \ell_2) + 1.
\end{align*}
\end{corollary}

\begin{theorem}
The spherical harmonic $Y_{\ell}(x)$ is an eigenvector of $W(x)$ for any $\SO$-equivariant $W : \R^3 \to \Hom(\R^{2\ell+1}, \R^{2\ell+1})$.
\end{theorem}

\begin{proof}
The $m = 0$ column of $D^{\ell}(g_x)$ is precisely $Y_{\ell}(\hat{x})$, up to normalization. This column spans the one-dimensional real-type invariant subspace $V_0$, on which $W(x)$ acts as a scalar.
\end{proof}


\section{The Low-Rank Algorithm}

The diagonalization theorem immediately yields an efficient algorithm for computing the matrix multiplication in an $\SO$-equivariant layer
\begin{align*}
	f : \left( \bigoplus_{\ell \in L} \R^{2\ell+1} \right)^m \to \left( \bigoplus_{\ell \in L'} \R^{2\ell+1} \right)^n.
\end{align*}

Write $p = \sum_{\ell \in L} (2\ell + 1)$ and $q = \sum_{\ell \in L'} (2\ell + 1)$ for the total dimensions.

\begin{enumerate}
	\item \textbf{Compute Wigner-D matrices.} Given $x \in \R^3$, compute $g_x \in \SO$ such that $g_x \cdot e_z = \hat{x}$. Form the block-diagonal matrices
	\begin{align*}
		P(x) = \bigoplus_{\ell \in L} D^{\ell}(g_x), \qquad Q(x) = \bigoplus_{\ell \in L'} D^{\ell}(g_x).
	\end{align*}
	These have shapes $p \times p$ and $q \times q$ respectively.

	\item \textbf{Transform to diagonal basis.} Multiply the input features $f \in \R^{p \times m}$ by $P(x)^{\top}$ to obtain $P(x)^{\top} f \in \R^{p \times m}$.

	\item \textbf{Apply radial weights.} The radial weights $\Lambda$ parameterize the block-diagonal structure. For each degree pair $(\ell_1, \ell_2)$:
	\begin{itemize}
		\item One scalar weight $\lambda_0$ for the real-type $m = 0$ block.
		\item Two weights $(a_m, b_m)$ for each complex-type $m > 0$ block, parameterizing
		$\begin{psmallmatrix} a_m & b_m \\ -b_m & a_m \end{psmallmatrix}$.
	\end{itemize}
	Apply these block-diagonal weights to obtain $\Lambda \, P(x)^{\top} f$.

	\item \textbf{Transform back.} Multiply by $Q(x)$ to obtain the output $Q(x) \Lambda \, P(x)^{\top} f \in \R^{q \times n}$.
\end{enumerate}

\subsection{Complexity Analysis}

The key advantage is that Step 3 operates on block-diagonal matrices rather than dense matrices. In the standard approach, one computes $\mathcal{O}(\lmax^3)$ basis matrices and performs tensor contractions of complexity $\mathcal{O}(\lmax^9)$.

In the low-rank approach:
\begin{itemize}
	\item Computing $P(x)$ and $Q(x)$ requires $\mathcal{O}(\lmax^2)$ Wigner-D matrix elements.
	\item The block-diagonal multiplication in Step 3 has complexity $\mathcal{O}(\lmax^2 \cdot mn)$.
	\item The dense multiplications in Steps 2 and 4 have complexity $\mathcal{O}(\lmax^4 \cdot m)$ and $\mathcal{O}(\lmax^4 \cdot n)$.
\end{itemize}
Overall complexity: $\mathcal{O}(\lmax^6)$ time and $\mathcal{O}(\lmax^4)$ space.


\subsection{Binned Radial Weights}

In molecular networks, the radial weights $\Lambda$ are typically parameterized by a neural network $\phi : \R_{\geq 0} \to \R^d$ that maps interatomic distances to weight vectors. For a batch of $B$ edges, the naive approach requires $B$ forward passes through $\phi$, and the resulting weight tensor of shape $(B, n, m, d)$ dominates memory usage during training.

\textbf{Observation.} The radial function $\phi$ is typically smooth and slowly varying. We can exploit this by discretizing the distance domain and using interpolation.

\textbf{Binned approach.} Partition $[0, r_{\max}]$ into $K$ bins with edges $0 = r_0 < r_1 < \cdots < r_K = r_{\max}$. Precompute a lookup table
\begin{align*}
	T_k = \phi(r_k) \in \R^{n \times m \times d}, \qquad k = 0, 1, \ldots, K.
\end{align*}
For an edge with distance $r \in [r_k, r_{k+1}]$, compute the interpolation weight $t = (r - r_k)/(r_{k+1} - r_k)$ and approximate
\begin{align*}
	\phi(r) \approx (1 - t) \, T_k + t \, T_{k+1}.
\end{align*}

\textbf{Complexity reduction.} This reduces MLP evaluations from $\mathcal{O}(B)$ to $\mathcal{O}(K)$, where typically $K \ll B$. More importantly, the weight tensor is never materialized: the fused CUDA kernel performs interpolation on-the-fly during the block-diagonal multiplication, reducing memory from $\mathcal{O}(B \cdot nmd)$ to $\mathcal{O}(K \cdot nmd)$.

\textbf{Gradient computation.} The interpolation is differentiable. During backpropagation, gradients with respect to the lookup table $T$ are computed via weighted scatter-add:
\begin{align*}
	\frac{\partial \mathcal{L}}{\partial T_k} = \sum_{i : r_i \in [r_{k-1}, r_k]} (1 - t_i) \frac{\partial \mathcal{L}}{\partial \phi(r_i)} + \sum_{i : r_i \in [r_k, r_{k+1}]} t_i \frac{\partial \mathcal{L}}{\partial \phi(r_i)}.
\end{align*}
This is implemented with atomic operations for efficiency. Gradients also flow through the interpolation weights $t$ to the distances $r$, enabling force computation.

\textbf{Accuracy.} For smooth radial functions, linear interpolation with $K = 100$ bins typically achieves relative error below $0.1\%$, which is negligible compared to other sources of approximation in neural network training.

\end{document}
